
Three research questions guided the experiments:

\textbf{RQ1:} Can automated DTS-weighted scoring achieve at least 70\% corpus coverage from public ML dataset repository APIs?

\textbf{RQ2:} Does DTS inverse-frequency weighting produce scores different from naive unweighted counting, and are the scores internally consistent?

\textbf{RQ3:} What is the per-section API coverage profile across the six DTS sections and three repositories?

\subsubsection{4.1 Corpus}

\begin{table}[htbp]
\centering
\resizebox{\columnwidth}{!}{%
\begin{tabular}{llll}
\toprule
Repository & Target N & Collected N & Stratification \\
\midrule
HuggingFace Hub & 500 & 496 & Task domain x upload year (8 bins) \\
OpenML & 200 & 200 & Task type \\
UCI ML Repository & ~100 & 62 & None (full accessible population) \\
**Total** & **~800** & **758** & -- \\
\bottomrule
\end{tabular}
}
\end{table}

\subsubsection{4.2 Comparison Reference}

This is a feasibility study. We compare against pre-specified gate thresholds (coverage >= 0.70, proxy r >= 0.70) and an unweighted scoring baseline (equal weight per field) for RQ2. Per-section rates are compared descriptively with the directional patterns reported in Rondina et al. (2025), though exact numerical comparison is not possible as the Rondina et al. values have not been independently verified.

\subsubsection{4.3 Evaluation Metrics}

\textbf{Coverage rate:} The proportion of datasets with a weighted DTS score greater than zero.

\textbf{Proxy Pearson r:} Correlation between weighted and unweighted DTS scores on a 120-dataset subsample, with 95\% bootstrap confidence interval (1,000 resamples).

\textbf{Per-section coverage rates:} The proportion of datasets in each repository with non-zero coverage for each DTS section.

\subsubsection{4.4 Implementation Details}

The pipeline is CPU-only. Collection required approximately 4 hours at the enforced rate limits. Scoring and validation completed in under 5 minutes. All random operations used seed 42.
